\section{Marco teórico}

\subsection{Fundamentos de series temporales}

Una serie temporal es un conjunto de observaciones registradas en instantes sucesivos, generalmente espaciados de manera uniforme. Su rasgo distintivo no es solo la secuencia cronológica, sino la dependencia entre observaciones vecinas: a diferencia de un conjunto transversal, una serie temporal incorpora memoria, persistencia, cambios de nivel, regularidades periódicas y perturbaciones irregulares \parencite{hyndman2021}. En el contexto de este proyecto, comprender esa estructura es indispensable porque la variable de interés ---la demanda eléctrica--- se observa secuencialmente y cualquier estrategia de preprocesamiento, descomposición o modelado depende de reconocer su naturaleza temporal.

Una forma clásica de interpretar una serie consiste en descomponerla en componentes. La tendencia representa la dirección de largo plazo sin implicar necesariamente crecimiento lineal; la estacionalidad alude a patrones que se repiten con periodicidad fija ---horas del día, días de la semana, estaciones---; los ciclos son oscilaciones de mediano o largo plazo sin periodicidad fija; y el ruido recoge la variación no explicada por los patrones sistemáticos \parencite{hyndman2021,harvey1989}. La autocorrelación complementa esta lectura al medir la relación lineal entre una observación y sus valores rezagados: la ACF resume cómo cambia esa dependencia conforme aumenta el lag, mostrando descensos lentos en presencia de tendencia y picos en múltiplos del periodo cuando hay estacionalidad \parencite{hyndman2021,box2015}. Esta memoria temporal ---el hecho de que el consumo a una hora se parezca al de la hora anterior o al de la misma hora del día previo--- es uno de los rasgos que vuelve razonable estudiar la serie desde su propia trayectoria histórica \parencite{hamilton1994}.

La estacionariedad, en su formulación débil, exige media constante, varianza constante y autocovarianza que dependa solo del lag. Cuando no se cumple ---por cambios de infraestructura, expansión de carga o nuevas rutinas de operación---, la diferenciación ayuda a estabilizar la media al remover cambios de nivel y tendencia, mientras que las transformaciones logarítmicas o Box-Cox ayudan a estabilizar la varianza \parencite{hyndman2021,box1964}. Conviene no confundir ambos recursos: la diferenciación actúa sobre la estructura dinámica y las transformaciones sobre la escala.

La descomposición temporal separa la serie en componentes interpretables. La forma aditiva (\(Y_t = T_t + S_t + E_t\)) es adecuada cuando la amplitud estacional es aproximadamente constante; la multiplicativa (\(Y_t = T_t \times S_t \times E_t\)), cuando dicha amplitud crece con el nivel \parencite{hyndman2021}. En demanda eléctrica, donde el aumento del consumo general puede amplificar las diferencias entre horas pico y valle, una descomposición puramente aditiva puede subrepresentar la dinámica real. El método STL ofrece una alternativa flexible al permitir estacionalidad cambiante y ser robusto a valores atípicos \parencite{hyndman2021}. Las transformaciones logarítmicas y Box-Cox complementan este análisis al estabilizar la varianza y convertir relaciones multiplicativas en aditivas, aunque su aplicación debe responder a la estructura empírica de la serie y no a una convención automática, ya que la retransformación puede afectar intervalos y sesgo \parencite{hyndman2021,box1964}.

\subsection{Fundamentos del pronóstico de series temporales}

El pronóstico de series temporales consiste en estimar valores futuros a partir de observaciones ordenadas cronológicamente. Cuatro decisiones metodológicas estrechamente conectadas definen este marco: el horizonte de predicción, la longitud de las ventanas de entrada, el procedimiento de validación temporal y las métricas de error \parencite{hyndman2021}.

El horizonte \(h\) indica cuántos pasos hacia adelante se desea estimar. No es equivalente pronosticar el siguiente punto que las próximas 24 horas: el horizonte define tanto la dificultad del problema como la utilidad operativa del resultado, y la incertidumbre tiende a incrementarse conforme \(h\) crece, especialmente en estrategias recursivas donde el error de un paso puede propagarse a los siguientes \parencite{tashman2000}. La ventana de entrada o lookback ---la cantidad de observaciones pasadas que se proporcionan al modelo--- condiciona de forma importante el desempeño: una ventana demasiado corta omite dependencias informativas como estacionalidad de corto alcance; una demasiado larga introduce redundancia, incrementa la dimensión de entrada y eleva el riesgo de sobreajuste \parencite{leites2024}. En series de demanda eléctrica, la selección del lookback suele guiarse por la frecuencia de muestreo, la periodicidad dominante y la estabilidad del régimen de consumo \parencite{hyndman2021,weron2006}.

La evaluación de modelos temporales no puede basarse en particiones aleatorias: estas rompen el orden cronológico y permiten que observaciones futuras influyan sobre predicciones del pasado \parencite{tashman2000}. La literatura recomienda procedimientos de rolling origin o validación cruzada temporal, donde el origen del pronóstico se desplaza repetidamente hacia adelante y se generan errores en múltiples orígenes, obteniendo una imagen más robusta del desempeño que una sola partición fija \parencite{hyndman2021}. Si bien Bergmeir, Hyndman y Koo mostraron que el k-fold estándar puede ser válido en condiciones autorregresivas específicas \parencite{bergmeir2018}, para forecasting operativo la recomendación dominante sigue favoreciendo la evaluación temporal porque refleja mejor el despliegue real.

En sistemas eléctricos, la predicción a corto plazo ---desde una hora hasta varios días--- es una tarea central para programación, seguridad y despacho \parencite{gross1987}. Su aparente facilidad puede ser engañosa: aunque la autocorrelación suele ser alta, también es el horizonte donde más pesan anomalías operativas, feriados y eventos atípicos \parencite{hippert2001}. En cuanto a las métricas, el MAE ofrece una lectura directa del desvío promedio en unidades reales; el RMSE penaliza con más fuerza los errores grandes, lo que lo hace pertinente cuando fallar en picos tiene un costo operativo alto; y el MAPE facilita la comunicación del error en términos porcentuales \parencite{hyndman2006}. Para este proyecto se propone usar conjuntamente MAPE, RMSE y MAE, ya que ofrecen perspectivas complementarias sobre precisión, severidad de errores extremos e interpretabilidad operativa.

\subsection{Demanda eléctrica como problema de modelado}

La demanda eléctrica es una serie temporal operativamente crítica y metodológicamente difícil porque combina autocorrelación, multiestacionalidad, sensibilidad al calendario, influencia meteorológica, eventos atípicos y cambios de régimen \parencite{alfares2002,hippert2001}. A diferencia de otras series más estables, la carga eléctrica presenta cambios sistemáticos por hora del día, día de la semana y estación del año, pero también responde a choques menos regulares como olas de calor, festivos móviles o interrupciones de red \parencite{alfares2002,taylor2008}.

La demanda está fuertemente estructurada por rutinas de uso. En la escala intradiaria, la forma de la carga depende de la alternancia entre periodos de baja actividad, arranque de actividades y horas de mayor ocupación; en la semanal, los días hábiles se diferencian de los fines de semana por intensidad y magnitud de los picos; en la anual, estaciones climáticas, calendarios escolares y periodos vacacionales modifican tanto la forma como el nivel general de consumo \parencite{alfares2002,hippert2001}. En series horarias suelen coexistir al menos tres ciclos ---diario, semanal y anual--- que no solo coexisten sino que interactúan: la forma de una hora determinada cambia según el día de la semana y la época del año \parencite{hyndman2021}. Esta multiestacionalidad explica por qué muchos métodos concebidos para una sola periodicidad dominante resultan insuficientes, y por qué se desarrollaron enfoques como TBATS para capturar patrones estacionales complejos \parencite{delivera2011}.

Aunque la demanda puede analizarse como serie endógena, la literatura coincide en que una parte sustantiva de su variabilidad proviene de variables externas. La temperatura es la más citada por su relación con enfriamiento y calefacción, y su impacto puede ser no lineal: German-Soto y Hernández Blanco muestran, para el caso mexicano, que el efecto del clima extremo se intensifica a partir de ciertos umbrales \parencite{germansoto2025}. La humedad puede complementar ese efecto en condiciones de calor húmedo \parencite{xie2018humidity}. Los festivos y días especiales alteran el patrón semanal regular, y tratarlos explícitamente ---distinguiendo entre festivos de fecha fija y asociados al día de la semana--- puede mejorar de forma sustantiva la precisión en los periodos más difíciles de pronosticar \parencite{ziel2018}.

La demanda también exhibe volatilidad local y episodios de pico que concentran presión sobre generación, transformación y reserva operativa. El costo de equivocarse durante un pico es mayor que en una hora plana, por lo que un modelo con buen error promedio puede seguir siendo inadecuado si falla sistemáticamente en las horas críticas \parencite{fan2012}. La complejidad operativa amplifica este problema: los operadores insertan el pronóstico en decisiones encadenadas sobre compromiso de unidades, reserva y gestión de picos, y mejorar la precisión del forecast se traduce en reducción de costos y menor necesidad de decisiones conservadoras \parencite{hobbs1999}. A ello se suma que no toda observación inusual representa el mismo fenómeno: algunas anomalías corresponden a errores de medición, otras a eventos reales del sistema y otras a cambios estructurales de comportamiento, y cada caso demanda un tratamiento distinto en el pipeline de pronóstico \parencite{luo2018}.

La literatura sobre demanda eléctrica resalta el valor de las variables exógenas. Sin embargo, el presente proyecto se delimita a un enfoque univariado con evaluación empírica. Esto implica modelar patrones internos de la propia serie ---tendencia, estacionalidad y recurrencias--- sin depender de información externa de temperatura, humedad o calendario codificado. Esta decisión es metodológicamente válida como primera solución aplicada con datos históricos disponibles, pero obliga a reconocer que en escenarios dominados por clima extremo o eventos atípicos la ausencia de covariables puede limitar la generalización.

\subsection{Modelos tradicionales de pronóstico}

Los modelos tradicionales de pronóstico permiten establecer referencias mínimas, interpretar con claridad la dinámica temporal y verificar si un modelo más complejo realmente agrega valor \parencite{hyndman2021,box2015}. En esta familia destacan los métodos naive, el suavizamiento exponencial y los modelos ARIMA/SARIMA \parencite{hyndman2021,degooijer2006}.

El método naive asume que el mejor pronóstico del siguiente periodo es el último valor observado; su versión estacional extiende la idea al último valor de la misma posición del ciclo. En demanda eléctrica, donde existen patrones diarios o semanales, esta referencia es pertinente porque captura de forma mínima la repetición de calendarios de consumo sin requerir ajuste sofisticado \parencite{hyndman2021,taylor2003}. La relevancia metodológica del naive se refuerza en la literatura de evaluación: el MASE, propuesto por Hyndman y Koehler, se escala respecto al error del método naive para permitir comparaciones interpretables entre series y modelos \parencite{hyndman2006}.

El suavizamiento exponencial asigna mayor peso a las observaciones más recientes. En las extensiones de Holt y Holt-Winters se incorporan tendencia y estacionalidad. La importancia teórica de esta familia se consolidó cuando Hyndman, Koehler, Snyder y Grose demostraron que estos métodos pueden representarse mediante modelos de espacio de estados, dando lugar al marco ETS (Error, Trend, Seasonal), que permite estimación por verosimilitud, cálculo de AIC y construcción de intervalos de predicción \parencite{hyndman2002}. Para demanda eléctrica, Taylor mostró que el suavizamiento exponencial con doble estacionalidad puede ser competitivo en pronóstico de corto plazo \parencite{taylor2003}, y la extensión a triple estacionalidad permitió capturar componentes diarios, semanales y anuales en datos de alta frecuencia \parencite{taylor2010}.

Los modelos ARIMA, formalizados por Box y Jenkins, modelan la serie a partir de la dependencia lineal respecto a sus rezagos y a los errores pasados, siempre que sea estacionaria o pueda volverse aproximadamente estacionaria mediante diferenciación \parencite{box2015}. La expresión ARIMA(p,d,q) resume tres componentes: \(p\) (orden autorregresivo), \(d\) (diferenciaciones) y \(q\) (orden de media móvil). Cuando la serie exhibe patrón estacional, ARIMA se amplía a SARIMA ---ARIMA(p,d,q)(P,D,Q)\(_s\)---, que incorpora términos estacionales con periodo \(s\) \parencite{hyndman2021,box2015}. Su limitación aparece cuando coexisten varias estacionalidades relevantes (intradía, intra-semana, intra-anual), lo que motivó el desarrollo de extensiones especializadas \parencite{taylor2003,taylor2010}.

La selección de modelos se apoya en criterios de información como AIC, AICc y BIC, que balancean bondad de ajuste y complejidad penalizando el número de parámetros \parencite{hyndman2021,hyndman2002}. Sin embargo, estos criterios deben complementarse con diagnóstico residual y validación temporal: un modelo no es adecuado solo porque minimiza AIC, sino que debe dejar residuos aproximadamente incorrelacionados y demostrar utilidad en particiones temporales de prueba \parencite{degooijer2006,chodakowska2021}.

La principal ventaja de los modelos tradicionales es su interpretabilidad estadística y su parsimonia: con pocas observaciones o ante estructuras predominantemente lineales, un modelo bien ajustado puede generalizar mejor que uno muy flexible \parencite{degooijer2006,hahn2009}. Su limitación fundamental aparece frente a no linealidades complejas, múltiples estacionalidades intrincadas o dependencia de variables exógenas omitidas. Sin embargo, de esta limitación no se sigue que carezcan de valor: metodológicamente constituyen el punto de referencia necesario para determinar si la mejora de un modelo más complejo es real o dependiente del escenario. Una vez establecidos estos referentes estadísticos, el siguiente paso es examinar arquitecturas de aprendizaje profundo que han mostrado capacidad para capturar patrones no lineales y dependencias complejas en series temporales.
